\section{Applications}

	\subsection{Series sums}
		\begin{enumerate}[\bf E.g. 1:]
			\item Write a Scilab function \verb|harmonic_sum(n)| that returns the sum
			\begin{align*}
				1+\dfrac{1}{2}+\dfrac{1}{3}+\cdots+\dfrac{1}{n}
			\end{align*}
			\textit{Solution:} We have to create a \code{.sci} file which can be later executed in the console window.
			\begin{enumerate}[I.]
				\item Contents of the \verb|harmonic_sum.sci| file:
				\lstinputlisting[language=Scilab]{code/harmonic_sum.sci}
				
				\item Execution on the console window:
				\begin{lstlisting}[language=Scilab]
					--> exec('harmonic_sum.sci', -1);
					--> harmonic_sum(15)
					ans = 
					3.3182290
				\end{lstlisting}
			\end{enumerate}
			
			\item Write a Scilab function \verb|newton_root(f, df, x0, tol, max_iter)| that implements the Newton-Raphson method to find a root of a function $f(x)$. Provide the function header and the essential loop (concise, runnable code), and briefly state two stopping criteria used in your implementation.
			
			\textit{Solution:} We have to create a \code{.sci} file which can be later executed in the console window.
			\begin{enumerate}[I.]
				\item Contents of the \verb|newton_root.sci| file:
				\lstinputlisting[language=Scilab]{code/newton_root.sci}
				
				\item Execution on the console window:
				\begin{lstlisting}[language=Scilab]
					--> exec('newton_root.sci', -1);
					--> deff('y=f1(x)', 'y = x^2 - 2');
					--> deff('y=df1(x)', 'y = 2*x');
					--> [root1, it1, conv1] = newton_root(f1, df1, 1.0, 1e-12, 50);
					--> mprintf('Test1: root = %.12f, iter = %d, converged = %d\n', root1, it1, conv1);
					Test1: root = 1.414213562373, iter = 6, converged = 1
				\end{lstlisting}
			\end{enumerate}
		\end{enumerate}